% !PRINTED 519
\weberpage{519}
Ist nun $m$ durch eine noch höhere als die zweite Potenz von 2 teilbar, so kann man, wenn der Exponent von 2 ungerade ist, nach § 105, (7) alle Charaktere auf solche von der Form
\[
\left(\frac{A}{m''}\right), \ \left(\frac{-1}{A}\right)\left(\frac{A}{m''}\right)
\]
zurückführen. Man setze
\[
m = 2r^2 m'',
\]
indem man wieder unter $r^2$ die größte in $\tfrac12 m$ aufgehende Quadratzahl versteht und wende, je nachdem $m'' \equiv 1$ oder $\equiv 3\ (\mathrm{mod}\ 4)$ ist, die Formel (8) oder (18) an. Man erhält dann die vollständige Zerfällung der Klassengleichung unter Adjunktion der Quadratwurzeln aus sämtlichen in $m$ aufgehenden Primzahlen, ausschließlich $\sqrt2$.

Ist aber der Exponent der höchsten in $m$ aufgehenden Potenz von 2 eine gerade Zahl, so genügt auch (18) noch nicht zur vollständigen Zerlegung der Klassengleichung. Man kann zwar hier wieder durch die Formeln (8), (18) durch Adjunktion sämtlicher $\sqrt p$ und $\sqrt q$ die Funktion $H$ zerfällen, die $\sqrt2$ bekommt man aber dadurch nicht hinein, und es bleiben also immer noch je zwei Geschlechter in einer solchen Teilgleichung enthalten. Um auch diese noch zu trennen, leiten wir unter der Voraussetzung, daß $m$ durch 8 teilbar sei, noch eine dritte Formel her.

Wir setzen
\[
m = m'm'',
\]
\[
x = 1, \quad y = \tfrac12 m' - 2m'', \quad n = \left(\tfrac14 m' + m''\right)^2.
\]
\[
\sqrt{-i(\alpha+\beta\omega)} = e^{-\frac{\pi i}{4}}\left(\tfrac12\sqrt{m'} + i\sqrt{m''}\right),
\]
worin wieder $m''$ ungerade und durch kein Quadrat teilbar ist. Hier folgt aus § 136:
\weq{19}{M = (-1)^{\frac{m}{8}} e^{\frac{\pi i}{4} A m''} \left(\frac{A}{m''}\right) \left(i^{-\frac{m''+1}{2}} \tfrac12\sqrt{m'} + i^{-\frac{m''-1}{2}}\sqrt{m''}\right)^3.}

Nachdem die Zerlegung nach den Charakteren $\left(\dfrac{A}{m''}\right)$ durch die Formeln (8) und (18) erledigt ist, handelt es sich nur noch um die Zerlegung nach den Charakteren
\[
\left(\frac{-1}{A}\right), \ \left(\frac{2}{A}\right),
\]
d. h. nach dem Verhalten von $A$ gegen den Modul 8. Es genügt daher, in der Formel (19) $m'' = 1$ zu setzen, wodurch man folgende vier Werte von $M$ erhält:
%CONT>>
